\section{Introduction}
\label{sec:introduction}

An autonomous navigation stack is architected to trust its perception
layer. Path planning, risk assessment and minimum-risk manoeuvre selection
all consume detections as though they were reliable, and the safety case
for the vehicle rests on that reliability holding under the conditions the
vehicle actually meets. Because it will not always hold, a monitoring layer
is attractive: something that watches perception during operation and tells
the planner when its inputs have stopped being trustworthy, so the stack
can slow down, hand back, or pull over.

Conformal prediction offers a principled way to build one. Calibrate a
nonconformity score on nominal data, and the fraction of new predictions
falling outside the conformal region is controlled at a chosen level
without distributional assumptions. Pair it with anytime-valid sequential
testing --- a betting process whose wealth is a nonnegative supermartingale
under the nominal hypothesis --- and the monitor can be interrogated
continuously without a multiple-testing penalty, which is exactly what a
vehicle needs from a system asked ``has anything changed yet?'' at every
frame.

We set out to reduce the detection delay of such a monitor for camera/LiDAR
3D detection under real adverse weather. We did not achieve that, and this
paper is about what we found instead. Two findings concern how monitors of
this design fail \emph{silently} --- continuing to emit calibrated
quantiles and plausible alarm times while not doing their job --- and both
matter more for navigation than a delay improvement would have.

\paragraph{A monitor can measure nothing and look healthy.}
The nonconformity score used throughout this literature is a residual
between predicted and ground-truth boxes. If the detector's regression head
degenerates --- as ours did, under a placeholder training loss whose
optimum is a constant --- that residual becomes $\lVert B_{\mathrm{target}}
\rVert_1$, a function of the labels alone. Every downstream component then
behaves exactly as specified: calibration returns a well-defined quantile,
because it is the quantile of \emph{some} exchangeable score distribution;
the betting process accumulates wealth, because the miscoverage rate still
moves; the alarm fires at a plausible time, because annotation statistics
genuinely differ between a dataset's nominal and adverse splits. The
monitor reports dataset structure as perception degradation, and nothing in
its output reveals this.

We introduce a \emph{null-detector diagnostic} that does reveal it: rerun
the whole pipeline, calibration included, with the prediction head's output
replaced by zeros. It costs one evaluation pass and needs no extra data,
training or annotation. Applied to our own earlier results it reproduced
them bit-identically --- a nominal-to-degraded jump of
$0.13904668284371138$ from both the real and the zeroed detector, agreeing
to 17 significant figures (Section~\ref{sec:null-diagnostic}). A monitor
with no detector performed identically to one with a detector, and the
reason is instructive rather than mysterious: content that is homogeneous
within a category produces a score that looks stable across sessions in a
way genuine, session-heterogeneous detector error does not, so a null
detector's output can in principle look \emph{no worse behaved} than the
real pipeline's even though it carries no perceptual information at all
(Section~\ref{sec:discuss-null}) --- precisely why the failure is invisible
from output inspection alone.

\paragraph{Exchangeability fails across driving sessions.}
Conformal validity is conditional on exchangeability between calibration
and test data. Real driving data arrives in sessions --- drives --- that
differ persistently in route, traffic, lighting and weather. On CADC
\cite{Pitropov2021CADC}, on a drive-disjoint held-out split, we measure
per-drive miscoverage under a single pooled quantile spreading $3$--$5\times$
across drives that are all nominal, all clear-road, on a nominal-drive
sample thin enough (four held-out drives) that we report the range as
indicative rather than precise (Section~\ref{sec:exchangeability}).
Between-session variation is comparable in size to the weather effect the
monitor exists to detect. The guarantee stays formally correct --- it holds
on average over the calibration distribution --- and becomes operationally
useless, because the monitored quantity is dominated by which session the
vehicle is on.

For navigation specifically this is the worst way for a monitor to fail.
An alarm that tracks route identity rather than perception loss will fire
on an unremarkable drive and stay silent on a dangerous one, and a planner
wired to react to it inherits that error directly. Worse, a monitor with
undemonstrated sensitivity that is nonetheless counted in a safety
argument is more dangerous than no monitor at all, because the false
assurance displaces other mitigations.

\paragraph{What we can and cannot say about the weather effect.}
Our original objective requires a detectable difference between nominal and
degraded conditions. On drive-disjoint held-out data, with resampling
clustered over both nominal sessions and degraded drives, we measure
$+0.069$ with a 95\% interval of $[-0.073, +0.210]$. This is an
\emph{underpowered null}, not a demonstrated absence, and we are careful to
report it as such throughout: the interval is consistent with a real effect
of the size we previously claimed and equally consistent with none. Nine
evaluation sessions cannot resolve an effect of this magnitude --- the
minimum detectable effect at 80\% power, computed directly from this
interval's bootstrap standard error, is roughly $0.20$, almost three times
the point estimate.

We also report why an earlier version of this measurement appeared
significant. Two defects compounded: the detector had trained on the very
drives being evaluated, and a bookkeeping error in the bootstrap let each
session's own calibration frames re-enter the window being monitored.
Reporting the interval both with and without the degraded arm resampled
shows that even the narrower comparison fails to reach significance once
the drive-disjoint split is applied (Section~\ref{sec:underpowered}). We
suggest dual reporting as a general practice for clustered effect
estimates: it helps a reader distinguish an underpowered design from a
genuine absence, which the point estimate alone never does.

\paragraph{Scope and honesty about this work's own history.}
Both findings above were discovered by controls applied to our own
pipeline, and the numbers those controls invalidated were ours. We think
this strengthens rather than weakens the case for the controls, but it puts
an obligation on us to be exact about what is being corrected: the affected
results were obtained in an unpublished predecessor of this work and were
never published, so no external record requires correction. Snowy
Scenes~\cite{SnowyScenes2026} appears throughout as a negative control: its
null detector shows a large, significant content-driven effect where its
real detector shows none, demonstrating that the diagnostic fires on a
genuine artifact when one is present. CADC's corrected result is more
modest than we first reported it to be --- neither its real nor its null
detector reaches significance --- which at least shows the diagnostic does
not fire spuriously when the honest answer is that we cannot tell.
